\exercisesection{FIELDS OF REPRESENTATIVES}

\begin{enumerate}
\def\labelenumi{\arabic{enumi})}
\tightlist
\item
  Let \(p\) be a prime number, \(G\) the free group on two generators \(X\) and \(Y\), and \(F_p[G]\) the group algebra of \(G\) over \(F_p\). Set \(Z = XY - YX\) and denote by \(\mathfrak{a}\) the two-sided ideal of \(F_p[G]\) generated by the elements \(ZX - XZ\), \(ZY - YZ\), and \(Z^2\). Put \(A = F_p[G]/\mathfrak{a}\), and denote by \(\mathfrak{z}\) the two-sided ideal of \(A\) generated by the image of \(Z\) in \(A\) (we will still denote by \(X\), \(Y\), \(Z\) the images in \(A\) of \(X\), \(Y\), and \(Z\)). a) Let \(a\) be an element of \(A\). Prove that there exist two finitely supported families of elements of \(F_p\), namely \((p_{\alpha,\beta})_{(\alpha,\beta)\in\mathbf{Z}^2}\) and \((q_{\alpha,\beta})_{(\alpha,\beta)\in\mathbf{Z}^2}\), uniquely determined by the condition
\end{enumerate}

\[
a = \sum_ {(\alpha , \beta) \in \mathbf {Z} ^ {2}} (p _ {\alpha , \beta} + q _ {\alpha , \beta} \mathbf {Z}) \mathrm{X} ^ {\alpha} \mathrm{Y} ^ {\beta}.
\]

Deduce that the quotient algebra \(A/\mathfrak{z}\) is commutative and integral, and that the ideal \(\mathfrak{z}\) is contained in the center of A.

\begin{enumerate}
\def\labelenumi{\alph{enumi})}
\setcounter{enumi}{1}
\tightlist
\item
  For every pair \((a, b)\) of elements of A, there exists an element c of A such that
\end{enumerate}

\[
  a ^ {k} b - b a ^ {k} = k a ^ {k - 1} c Z \quad \text { for every integer } k \geqslant 0.
\]

From this, deduce that the set \(A^{p}\) of the p-th powers of the elements of A is contained in the center of A.

\begin{enumerate}
\def\labelenumi{\alph{enumi})}
\setcounter{enumi}{2}
\tightlist
\item
  The multiplicative part \(S = A - \mathfrak{z}\) of \(A\) allows a computation of fractions on the right and on the left (II, § 2, exerc. 22). We denote by B the ring of fractions of \(A\) thus constructed.
\end{enumerate}

(To verify, for example, that for every \(a\) in A and every \(s\) in S, there exists \(b \in A\) and \(t \in S\) such that \(at = sb\), take \(b = s^{p-1}a\), \(t = s^p\), and use b).)

\begin{enumerate}
\def\labelenumi{\alph{enumi})}
\setcounter{enumi}{3}
\item
  Prove that the canonical map from A to B is injective. (Note that the kernel of this map is contained in \(\mathfrak{z}\), and reason as in (a)).
\item
  The center of B contains Z. If m is the two-sided ideal of B generated by Z, then \(m^{2} = 0\).
\item
  The ideal m is the set of non-invertible elements of B, and the field B/m is isomorphic to the field \(\mathbf{F}_{p}(\mathbf{X}, \mathbf{Y})\) .
\item
  Prove that there is no subfield of B whose image in B/m is all of B/m.
\end{enumerate}

In Exercises 2 to 5, the (commutative) topological rings are assumed to be linearly topologized. If A and B are two topological rings and B is an A-algebra, one says that B is a topological A-algebra if the homomorphism from A to B defining the algebra structure is continuous.

\begin{enumerate}
\def\labelenumi{\arabic{enumi})}
\setcounter{enumi}{1}
\tightlist
\item
  Let A be a topological ring, and let B be a topological A-algebra. We say that B is a formally smooth A-algebra \(^{1}\) if, for every discrete topological A-algebra C, every ideal j of C with square zero, and every continuous A-homomorphism \(u: B \to C/j\) there exists a continuous A-homomorphism \(v: B \to C\) which, by passing to the quotient, induces u.
\end{enumerate}

Let A be a ring, and B an A-algebra. We say that B is a smooth A-algebra if B is a formally smooth A-algebra when A and B are endowed with the discrete topologies.

\begin{enumerate}
\def\labelenumi{\alph{enumi})}
\item
  Let A be a topological ring and B a formally smooth A-algebra. Let C be a ring and j an ideal of C. Endow C with the j-adic topology, and suppose C is separated and complete for this topology. Then, for every continuous A-homomorphism \(u:B\to C/\mathfrak{j}\), there exists a continuous A-homomorphism \(v:B\to C\) which, by passing to the quotient, induces \(u\).
\item
  Let A be a ring. Every polynomial algebra over A is a smooth A-algebra.
\item
  Let A be a topological ring. If B is a formally smooth A-algebra and C is a formally smooth B-algebra, then C is a formally smooth A-algebra.
\item
  Let A be a topological ring, B a formally smooth A-algebra, and A' a topological A-algebra. Then the topological A'-algebra \(B \otimes_{A} A'\) (III, § 2, exerc. 28) is a formally smooth A'-algebra.
\item
  Let A be a topological ring, and let B be a formally smooth A-algebra. Let S (resp. T) be a multiplicative subset of A (resp. B) such that the image of S in B is contained in T. Then \(T^{-1}B\) is a formally smooth \(S^{-1}A\)-algebra (see III, § 2, exerc. 27 for the topology of \(T^{-1}B\) and \(S^{-1}A\)).
\item
  Let n be an integer \(\geqslant 1\), and let \((\mathbf{B}_{i})_{1\leqslant i\leqslant n}\) be a family of topological A-algebras. For \(\prod_{i=1}^{n} B_{i}\) to be a formally smooth A-algebra, it is necessary and sufficient that \(B_{i}\) be a formally smooth A-algebra for \(1 \leqslant i \leqslant n\).
\item
  Let A be a topological ring, B a topological A-algebra, \(\hat{A}\) and \(\hat{B}\) the respective separated completions of A and B. Then the following three conditions are equivalent:
\end{enumerate}

\begin{enumerate}
\def\labelenumi{(\roman{enumi})}
\item
  B is a formally smooth A-algebra;
\item
  \(\hat{\mathbf{B}}\) is a formally smooth A-algebra;
\item
  \(\hat{\mathbf{B}}\) is a formally smooth \(\hat{\mathbf{A}}\)-algebra.
\end{enumerate}

\(h)\) We resume the notation and hypotheses of \(e\)). Then \(\mathbf{B}\{\mathrm{T}^{-1}\}\) is a formally smooth \(\mathbf{A}\{\mathrm{S}^{-1}\}\)-algebra.

\begin{enumerate}
\def\labelenumi{\arabic{enumi})}
\setcounter{enumi}{2}
\tightlist
\item
  Let \(k\) be a field and A a commutative \(k\)-algebra. For every integer \(i \geqslant 1\), set \(\mathbf{B}_i = \underbrace{\mathbf{A} \otimes_k \mathbf{A} \otimes_k \cdots \otimes_k \mathbf{A}}_{i \text{ terms}}\) and endow \(\mathbf{B}_i\) with the structure of an A-algebra obtained by letting
\end{enumerate}

A acts on the first component. Set \(C_{1} = B_{2}\), \(C_{2} = B_{3}\), \(C_{3} = B_{4} \oplus B_{3}\), and define A-linear maps \(d_{2}:C_{2} \to C_{1}\) and \(d_{3}:C_{3} \to C_{2}\) by the formulas

\[
d _ {3} ((1 \otimes x \otimes y \otimes z, 1 \otimes \alpha \otimes \beta)) = x \otimes y \otimes z - 1 \otimes x y \otimes z + 1 \otimes x \otimes y z - z \otimes x \otimes y + 1 \otimes \alpha \otimes \beta - 1 \otimes \beta \otimes \alpha,
\]

\[
d _ {2} (1 \otimes \alpha \otimes \beta) = \alpha \otimes \beta - 1 \otimes \alpha \beta + \beta \otimes \alpha ,
\]

\(^{1}\) If A is a local, Noetherian, complete algebra over a (discrete) field k, this definition coincides with that given in VIII, p.~98, exerc. 30, according to exerc. 5 below.

for any choice of elements \(x, y, z, \alpha, \beta\) of A. One thus obtains a complex of A-modules

\[
\mathrm{C} _ {k} (\mathrm{A}): \mathrm{C} _ {3} \xrightarrow {d _ {3}} \mathrm{C} _ {2} \xrightarrow {d _ {2}} \mathrm{C} _ {1}.
\]

If N is an A-module, one says that a \(k\)-bilinear map \(f: \mathbf{A} \times \mathbf{A} \to \mathbf{N}\) is a 2-cocycle if one has the identity \(xf(y, z) - f(xy, z) + f(x, yz) - zf(x, y) = 0\) for \(x, y, z\) in A, and that it is symmetric if one has \(f(x, y) = f(y, x)\) for every pair \((x, y) \in \mathbf{A} \times \mathbf{A}\).

\begin{enumerate}
\def\labelenumi{\alph{enumi})}
\tightlist
\item
  The following conditions are equivalent:
\end{enumerate}

\begin{enumerate}
\def\labelenumi{(\roman{enumi})}
\item
  the complex \(\operatorname{Hom}_{\mathbf{A}}(\mathbf{C}_k(\mathbf{A}), \mathbf{N})\) is acyclic for every A-module \(\mathbf{N}\) ;
\item
  for every A-module N, every symmetric 2-cocycle \(f:A\times A\to\mathbf{N}\) is a 1-coboundary, that is, of the form \(f(a,b)=ag(b)+bg(a)-g(ab)\) with \(g\in\mathrm{Hom}_{k}(\mathbf{A},\mathbf{N})\);
\item
  A is a smooth algebra over k.
\end{enumerate}

\begin{enumerate}
\def\labelenumi{\alph{enumi})}
\setcounter{enumi}{1}
\tightlist
\item
  If A is a field, the preceding conditions are also equivalent to the condition that the complex \(C_{k}(A)\) be acyclic.
\end{enumerate}

\begin{enumerate}
\def\labelenumi{\arabic{enumi})}
\setcounter{enumi}{3}
\tightlist
\item
  \begin{enumerate}
  \def\labelenumii{\alph{enumii})}
  \tightlist
  \item
    Let \(k\) be a field and \(K\) an extension of \(k\). If \(K\) is separable over \(k\), then \(K\) is a smooth \(k\)-algebra. (If \(K\) is a finitely generated extension of \(k\), use prop. 1 of § 3, n° 2. In the general case, write \(K\) as a union of finitely generated extensions of \(k\) and use the preceding exercise.)
  \end{enumerate}
\end{enumerate}

\begin{enumerate}
\def\labelenumi{\alph{enumi})}
\setcounter{enumi}{1}
\tightlist
\item
  Let A be a local ring that is separated and complete, containing a field k. Using a), prove that A possesses a field of representatives. If the residue field of A is a separable extension of k, then A possesses a field of representatives containing k.
\end{enumerate}

\begin{enumerate}
\def\labelenumi{\arabic{enumi})}
\setcounter{enumi}{4}
\tightlist
\item
  Let A be a Noetherian local ring containing a field \(k\) and m the maximal ideal of A. We assume that A is a formally smooth \(k\)-algebra (the field \(k\) being discrete).
\end{enumerate}

\begin{enumerate}
\def\labelenumi{\alph{enumi})}
\item
  Let K be a field of representatives of the ring \(\mathbb{A}/\mathfrak{m}^2\) (exercise 4, b)), and let \(x_1, \ldots, x_d\) be elements of m whose images form a basis of the K-vector space \(\mathfrak{m}/\mathfrak{m}^2\). Let \(\mathrm{K}[\mathrm{X}_1, \ldots, \mathrm{X}_d]\) be the ring of polynomials in \(d\) variables \(\mathrm{X}_1, \ldots, \mathrm{X}_d\), let \(\mathfrak{n}\) be its ideal generated by \(\mathrm{X}_1, \ldots, \mathrm{X}_d\), and let \(\varphi\) be the homomorphism from \(\mathrm{K}[\mathrm{X}_1, \ldots, \mathrm{X}_d]\) into \(\mathbb{A}/\mathfrak{m}^2\) which sends \(\mathrm{X}_i\) to \(x_i\) for \(1 \leqslant i \leqslant d\) and induces the identity on K. Then \(\varphi\) induces an isomorphism \(\overline{\varphi}\) from \(\mathrm{K}[\mathrm{X}_1, \ldots, \mathrm{X}_d]/\mathfrak{n}^2\) onto \(\mathbb{A}/\mathfrak{m}^2\).
\item
  For every integer \(n \geqslant 1\), there exists a homomorphism \(\psi_n: \mathbb{A} \to \mathrm{K}[\mathrm{X}_1, \ldots, \mathrm{X}_d]/\mathfrak{n}^{n+1}\) which, by passage to quotients, induces \(\overline{\varphi}^{-1}\). Such a homomorphism \(\psi_n\) is surjective.
\item
  For every integer \(n \geqslant 1\), the length of the A-module \(\mathbf{A} / \mathfrak{m}^{n+1}\) is at least \(\binom{d+n}{d}\), and one has \(\dim(\mathbf{A}) \geqslant d\).
\item
  For every finite extension \(k'\) of \(k\), and every maximal ideal \(m'\) of the semi-local ring \(A \otimes_k k'\), the local ring \((A \otimes_k k')_{m'}\) is regular.
\end{enumerate}

\begin{enumerate}
\def\labelenumi{\arabic{enumi})}
\setcounter{enumi}{5}
\tightlist
\item
  Let \(k\) be a field and \(K\) an extension of \(k\). Assume that \(K\) is a smooth \(k\)-algebra. Let \(k'\) be an algebraic extension of finite degree of \(k\).
\end{enumerate}

\begin{enumerate}
\def\labelenumi{\alph{enumi})}
\item
  The ring K \(\otimes_{k} k'\) is a product of a finite number of Artinian local rings, which are smooth \(k'\)-algebras (use exerc. 2, p.~76).
\item
  The ring K \(\otimes_{k} k'\) is reduced (use the preceding exercise).
\item
  The field K is separable over k.
\end{enumerate}

\begin{enumerate}
\def\labelenumi{\arabic{enumi})}
\setcounter{enumi}{6}
\tightlist
\item
  Let \(k\) be a field, \(P\) its prime subfield, and \(K\) an extension of \(k\). Then the following conditions are equivalent:
\end{enumerate}

\begin{enumerate}
\def\labelenumi{\alph{enumi})}
\item
  Every derivation of k into a K-module M extends uniquely to a derivation of K into M.
\item
  We have \(\Omega_{\mathbf{P}}(\mathbf{K}) = \Omega_{\mathbf{P}}(k)\otimes_k\mathbf{K}\).
\item
  The field K is separable over k, and \(\Omega_{k}(\mathbf{K}) = 0\) .
\end{enumerate}

If \(k\) is of characteristic 0, these conditions are also equivalent to the condition:

\begin{enumerate}
\def\labelenumi{\alph{enumi})}
\setcounter{enumi}{3}
\tightlist
\item
  The field K is an algebraic extension of k.
\end{enumerate}

If \(k\) has characteristic \(p\ne0\), they are also equivalent to the conditions:

\begin{enumerate}
\def\labelenumi{\alph{enumi})}
\setcounter{enumi}{4}
\item
  \(\mathbf{K} = k\otimes_{k^p}\mathbf{K}^p.\)
\item
  Every p-basis of k is also a p-basis of K.
\end{enumerate}

\begin{enumerate}
\def\labelenumi{\arabic{enumi})}
\setcounter{enumi}{7}
\tightlist
\item
  Let \(k\) be a field and \(K\) an extension of \(k\). We say that \(K\) is formally étale over \(k\) if, for every \(k\)-algebra \(A\), every ideal \(a\) of \(A\) with square zero, and every \(k\)-homomorphism \(u:K\to A/a\), there exists a unique \(k\)-homomorphism \(v:K\to A\) which gives \(u\) by passing to the quotient.
\end{enumerate}

\begin{enumerate}
\def\labelenumi{\alph{enumi})}
\item
  If K is formally étale over k, the equivalent conditions of the preceding exercise are satisfied. (If M is a K-module, consider the k-algebra whose underlying k-vector space is K ⊕ M, with M being an ideal of square zero.)
\item
  Conversely, if the equivalent conditions of the preceding exercise are satisfied, then K is formally étale over k. (Since the field K is separable over k, Exercise 4 allows one to prove the existence of v.)
\item
  Let \(\mathbf{B} = (b_{i})_{i\in\mathbf{I}}\) be a family of elements of K such that \((db_{i})_{i\in\mathbf{I}}\) is a basis of \(\Omega_{k}(\mathbf{K})\) over K. If K is separable over k, then \(k(\mathbf{B})\) is a purely transcendental extension of k (use Th. 2 of A, V, p.~125, Th. 1 of A, V, p.~97, and Exerc. 6 of A, V, p.~165) and K is formally étale over \(k(\mathbf{B})\).
\end{enumerate}

\begin{enumerate}
\def\labelenumi{\arabic{enumi})}
\setcounter{enumi}{8}
\tightlist
\item
  Let A be a local ring of equal characteristics, \(m_{A}\) its maximal ideal, and k a subfield of A. Then the residue field \(\kappa_{A}\) is a k-algebra. One says that k is a weak residue field of A if \(\kappa_{A}\) is formally étale over k (exerc. 8).
\end{enumerate}

\begin{enumerate}
\def\labelenumi{\alph{enumi})}
\item
  If A has a field of weak representatives \(k\), then the separated completion \(\hat{\mathbf{A}}\) of A has a unique field of representatives containing the image of \(k\) in \(\hat{\mathbf{A}}\).
\item
  Let B be a local ring contained in A, with maximal ideal \(m_{B} = m_{A} \cap B\). If \(\kappa_{A}\) is a separable extension of the residue field \(\kappa_{B}\) of B, every weak residue field of B is contained in a weak residue field of A (use Exercise 8(c)).
\item
  Let B be a local ring contained in A, with maximal ideal \(m_{B} = m_{A} \cap B\). Suppose further that A has characteristic \(p\ne0\) and contains \(B^{p}\). Then there exists a weak field of representatives of A which contains a weak field of representatives of B.
\end{enumerate}
% semantic-fragments: legacy-input-seam
\relax{}
